\section{Running time and concluding remarks}\label{sec:v4i-conclusion}

\subsection{Proof of Theorem~\ref{thm:main}}
Embeddings are output combinatorially, as rotation systems together with
the relative positions of the components of each colour. A pair of them is
\emph{compatible} if it comes from a simultaneous embedding.

Normalize in $O(N)$ time (Theorem~\ref{thm:v4-normalize}); the factors have
total size $O(N)$. A factor whose common graph has no edges is handled
directly, since any pair of planar rotation systems is compatible. For a
factor of size $n$ with a common cycle, construct the compiled instance $I$
in $O(n^2)$ time. It has size $O(n)$ (Theorem~\ref{thm:v4-compile}), so in
Theorem~\ref{thm:eng4-main} the initial size $S_0$, the number $m_0$ of
edges and the maximum pipe degree $D_0$ are $O(n)$, and the engine decides
$I$ and constructs a witness in $O\bigl((m_0+1)(S_0+m_0D_0)\bigr)=O(n^3)$
time. By Theorem~\ref{thm:v4-compile} the factor has a simultaneous
embedding exactly when $I$ is feasible, and a witness of $I$ is decoded and
repaired into a compatible pair for the factor in $O(n^2)$ time
(Theorem~\ref{thm:v4-repair}). Since $\sum_t n_t^3\le(\sum_t n_t)^3$, all
factors together take $O(N^3)$ time, and the normalization is inverted in
$O(N)$ time (Theorem~\ref{thm:v4-normalize}). \qed

\subsection{Common components that are cacti or biconnected}
Call a component of $H$ that is $2$-connected with at least three vertices a
\emph{$2$-component}.

\begin{corollary}\label{cor:v4i-mixed}
\sefe{} can be decided in $O(N^3)$ time for instances in which every
connected component of the common graph is a cactus or $2$-connected. For a
positive instance a simultaneous embedding is constructed in polynomial
time.
\end{corollary}

\begin{proof}
We apply the preprocessing of~\cite[Section~4]{BKR18} to the instance as
given, before the normalization. Let $C_1,\dots,C_k$ be the
$2$-components. In linear time it either finds that some $C_i$ has no
admissible embedding, in which case we reject, or computes admissible
embeddings of all $C_i$ and instances $I_1,\dots,I_m$ of total size $O(N)$,
one for each union subbridge with respect to the $C_i$, such that the input
has a simultaneous embedding if and only if every $I_j$ has
one~\cite[Lemmas~7 and~9]{BKR18}. Lemma~9 there is stated for connected
inputs, but its proof does not use connectivity. In $I_j$ every $C_i$ with
an attachment in the subbridge is replaced by a cycle through these
attachments in the order of the admissible embedding, with a fresh vertex
inserted between any two consecutive attachments (two fresh vertices if
there is only one attachment); a subbridge without attachments is a
component of $U$ and is kept unchanged. The new cycle subdivides the cycle
of~\cite{BKR18} (with a triangle at the attachment when there is only one,
which can be drawn in a corner there), and subdividing common edges
preserves the existence of a simultaneous embedding, as in the proof
of~\cite[Lemma~9]{BKR18}. Every new edge has a fresh end, so the
inputs stay simple. If an $I_j$ has a nonplanar colour we reject.

Let $K$ be a component of $H$ other than the $C_i$. If an edge $e$ of $K$
lies in a subbridge $S$, then for every $i$ both $S-V(C_i)$ and $K=K-V(C_i)$
are connected and share the ends of $e$, so maximality of $S$ gives
$K\subseteq S$; an isolated vertex of $H$ likewise lies in a single
subbridge. No common edge joins $K$ to a $C_i$, and the replacement cycles
use only vertices of the $C_i$ and fresh vertices. Hence the common graph of
every $I_j$ is a disjoint union of components of $H$, which are cacti, and
of cycles, so it is a cactus. Theorem~\ref{thm:main} decides each $I_j$ in
$O(N_j^3)$ time, and $\sum_jN_j^3\le(\sum_jN_j)^3=O(N^3)$. For a positive
instance, the admissible embeddings of the $C_i$ and the simultaneous
embeddings of the $I_j$ are glued face by face as in the proofs
of~\cite[Lemmas~7 and~9]{BKR18}.
\end{proof}

The reduction must precede the normalization of
Theorem~\ref{thm:v4-normalize}, whose augmentation may attach common edges
to a $2$-component.

\subsection{Limits and open questions}
The one feature of our engine without a counterpart in plain Synchronized
Planarity is the observation of a gap at a P-vertex, a bit that is a
function of an unrestricted cyclic order. Deleting a self-pipe while keeping
only the equalities between observations turns an infeasible instance into
a feasible one (Remark~\ref{rem:eng4-selfpipe}), and arbitrary families of
gap observations at one vertex need not be describable by parity
constraints (Remark~\ref{rem:eng4-anchors}). Neither example excludes other
gadgets, and a reduction of the compiled instances to plain Synchronized
Planarity remains open. Schaefer conjectured a strong Hanani--Tutte
theorem for \sefe{}~\cite[Conjecture~6.20]{Sch13}, which would have given a
parity-based decision algorithm; the conjecture is false, since Gutwenger,
Mutzel and Schaefer~\cite{GMS14} turned the clustered-planarity
counterexamples of Fulek, Kyn\v{c}l, Malinovi\'c and
P\'alv\"olgyi~\cite[p.~6]{FKMP15} into counterexamples to it; see
also~\cite[p.~137]{FPS17}.

The complexity of \sefe{} itself remains open. Since \sefe{} is polynomial
for connected common graphs~\cite{BFR21}, Corollary~\ref{cor:v4i-mixed}
implies that, unless P${}={}$NP, an NP-hardness reduction must produce, for
infinitely many inputs, disconnected common graphs with a component that
contains a block that is neither an edge nor a cycle and shares a cutvertex
with a further block. An
extension of our method would have to handle such blocks: their parallel
skeletons can display several common cycles, so the gap observations at
one vertex refer to several anchor pairs, and the cycles through a vertex
no longer form a flower.

We have not tried to optimize the running time. The compiled instances and
the terminal solver are linear, and Synchronized Planarity runs in quadratic
time~\cite{BFR21}. Our cubic bound has two sources. First, we keep the
leaves of the trees inserted by PropagatePQ as degree-two vertices, which
the size bounds of~\cite[Lemma~10(ii),(iii) and Theorem~12]{BFR21} do not
cover, so we bound all sizes by $O(S_0+m_0D_0)$ only. Second, every
iteration recomputes the wheel graphs, planarity, blocks, SPQR trees and the
scheduling data, and runs the conversion and the backward maps, on the
whole instance. We expect that both can be removed. A fresh leaf carries no
observation, pipe end or Q-cell variable before ConvertSmall, so it can be
suppressed at that point, which restores the leaf convention
of~\cite{BFR21}; and records kept per connected component, rebuilt only for
the components an operation changes, together with append-only parity rows,
should make every iteration cost $O(m_0)$. This would give $O(S_0+m_0^2)$
for Theorem~\ref{thm:eng4-main} and $O(N^2)$ for Theorem~\ref{thm:main}; we
have not carried out the details.
